
\section*{\SGAInlineRefTarget{sga3:i:section:5}5. Group cohomology}
\addcontentsline{toc}{section}{5. Group cohomology}

\subsection*{\SGAInlineRefTarget{sga3:i:subsection:5-1}5.1. The standard complex}

Let \(\mathscr C\) be a category, \(\mathbf G\) a
\(\widehat{\mathscr C}\)-group, \(\mathbf O\) a
\(\widehat{\mathscr C}\)-ring, and \(\mathbf F\) a
\(\mathbf G\)-\(\mathbf O\)-module. Set
\[
  C^n(\mathbf G,\mathbf F)
    =\operatorname{Hom}(\mathbf G^n,\mathbf F),
  \qquad
  \underline C^{\,n}(\mathbf G,\mathbf F)
    =\underline{\operatorname{Hom}}(\mathbf G^n,\mathbf F),
  \qquad n\geq0,
\]
where \(\mathbf G^0\) is the final object \(\underline{\mathbf e}\).
Then \(\underline C^{\,n}(\mathbf G,\mathbf F)\) (respectively,
\(C^n(\mathbf G,\mathbf F)\)) is endowed in the evident way with an
\(\mathbf O\)-module structure (respectively, a
\(\Gamma(\mathbf O)\)-module structure), and
\[
  C^n(\mathbf G,\mathbf F)
    \simeq \Gamma\bigl(\underline C^{\,n}(\mathbf G,\mathbf F)\bigr),
  \qquad
  \underline C^{\,n}(\mathbf G,\mathbf F)(S)
    =C^n(\mathbf G_S,\mathbf F_S).
\]

To give an element of \(C^n(\mathbf G,\mathbf F)\) is to give, for every
\(S\in\operatorname{Ob}(\mathscr C)\), an \(n\)-cochain of
\(\mathbf G(S)\) with values in \(\mathbf F(S)\), functorially in \(S\).
Recall that the boundary operator
\[
  \partial:C^n\bigl(\mathbf G(S),\mathbf F(S)\bigr)
    \longrightarrow
    C^{n+1}\bigl(\mathbf G(S),\mathbf F(S)\bigr)
\]
is given by\footnote{Translator's note: in the middle sum the 1970 text has
\((-1)^nf(g_1,\ldots,g_ig_{i+1},\ldots,g_n)\), misprints for
\((-1)^if(g_1,\ldots,g_ig_{i+1},\ldots,g_{n+1})\), which the Gille--Polo
re-edition (2011), Exposé~I, gives.}
\begin{align*}
  \partial f(g_1,\ldots,g_{n+1})
    ={}&g_1f(g_2,\ldots,g_{n+1}) \\
     &+\sum_{i=1}^{n}(-1)^i
       f(g_1,\ldots,g_i g_{i+1},\ldots,g_{n+1}) \\
     &+(-1)^{n+1}f(g_1,\ldots,g_n).
\end{align*}
It is functorial in \(S\), and hence defines a homomorphism
\[
  \partial:C^n(\mathbf G,\mathbf F)
    \longrightarrow C^{n+1}(\mathbf G,\mathbf F)
\]
such that \(\partial\circ\partial=0\).\footnote{Translator's note: the Gille--Polo re-edition (2011), Exposé~I, also writes
\(\partial\) without elements. In our formulation: if
\(m:\mathbf G\times\mathbf G\to\mathbf G\) is the multiplication and
\(\mu:\mathbf G\times\mathbf F\to\mathbf F\) the action, then for
\(f:\mathbf G^n\to\mathbf F\) the morphism
\(\partial f:\mathbf G^{n+1}\to\mathbf F\) is the alternating sum, with signs
\((-1)^0,\ldots,(-1)^{n+1}\), of \(\mu\circ(\operatorname{id}_{\mathbf G}\times f)\),
of the \(n\) composites of \(f\) with the morphisms
\(\mathbf G^{n+1}\to\mathbf G^n\) multiplying two adjacent factors, and of
the composite of \(f\) with the projection forgetting the last factor; the
same formula, after base change, gives \(\partial\) on
\(\underline C^{\,n}(\mathbf G,\mathbf F)(S)=C^n(\mathbf G_S,\mathbf F_S)\).}
We have thus defined a complex of
abelian groups (indeed, of \(\Gamma(\mathbf O)\)-modules), denoted
\(C^*(\mathbf G,\mathbf F)\).\footnote{Translator's note: this complex is often called the Hochschild
complex; see M.~Demazure and P.~Gabriel, \textit{Groupes algébriques},
tome~I, Masson and North-Holland, 1970, II, \S3 (a reference given in the
Gille--Polo re-edition (2011), Exposé~I).} In the same way one defines the complex of
\(\mathbf O\)-modules \(\underline C^{\,*}(\mathbf G,\mathbf F)\), and
\[
  C^*(\mathbf G,\mathbf F)
    =\Gamma\bigl(\underline C^{\,*}(\mathbf G,\mathbf F)\bigr).
\]
We denote by \(H^n(\mathbf G,\mathbf F)\) (respectively,
\(\underline H^{\,n}(\mathbf G,\mathbf F)\)) the groups (respectively,
the \(\widehat{\mathscr C}\)-groups) of cohomology of
\(C^*(\mathbf G,\mathbf F)\) (respectively,
\(\underline C^{\,*}(\mathbf G,\mathbf F)\)). In particular,
\[
  \underline H^{\,0}(\mathbf G,\mathbf F)=\mathbf F^{\mathbf G},
  \qquad
  H^0(\mathbf G,\mathbf F)=\Gamma(\mathbf F^{\mathbf G}).
\]

\subsection*{\SGAInlineRefTarget{sga3:i:subsection:5-2}5.2.}

We shall now show that
the functors \(\underline H^{\,n}\) (respectively, \(H^n\)) are indeed the
derived functors of \(\underline H^{\,0}\) (respectively, \(H^0\)).\footnote{Translator's note: the 1970 text has no heading 5.2 here, although it
numbers 5.2.1 and 5.2.2 and refers to 5.2 in the proof of 5.3.1; the heading
is supplied in the Gille--Polo re-edition (2011), Exposé~I.}

Let \(\mathbf P\) be an \(\mathbf O\)-module, and consider the object
\(\underline{\operatorname{Hom}}(\mathbf G,\mathbf P)\) of
\(\widehat{\mathscr C}\). It can be given a
\(\mathbf G\)-\(\mathbf O\)-module structure by letting \(\mathbf G\) act on
the first factor and \(\mathbf O\) on the second. Precisely,
\(\underline{\operatorname{Hom}}(\mathbf G,\mathbf P)(S)
=\operatorname{Hom}_S(\mathbf G_S,\mathbf P_S)\), and \(g\in\mathbf G(S)\)
and \(a\in\mathbf O(S)\) act by the formulas
\[
  (gf)(x)=f(xg),
  \qquad
  (af)(x)=af(x),
  \qquad x\in\mathbf G(S').
\]
Now let \(\mathbf F\) be a \(\mathbf G\)-\(\mathbf O\)-module, and let
\(E(\mathbf F)\) denote the object
\(\underline{\operatorname{Hom}}(\mathbf G,\mathbf F)\) with the preceding
\(\mathbf G\)-\(\mathbf O\)-module structure. There is a morphism of
\(\mathbf G\)-\(\mathbf O\)-modules
\[
  \mu_{\mathbf F}:\mathbf F\longrightarrow E(\mathbf F),
\]
defined as follows: with every arrow \(S'\to S''\) of \(\mathscr C\) one
must associate, functorially, a map from \(\mathbf F(S')\) to the set
\(\operatorname{Hom}(\mathbf G(S'),\mathbf F(S'))\); for this, associate
with every \(f\in\mathbf F(S')\) the map \(k_f\) defined by
\(k_f(g)=gf\). One checks at once that \(\mu_{\mathbf F}\) is compatible
with the \(\mathbf G\)-\(\mathbf O\)-module structures, and that it is even
a monomorphism.

\par\medskip\noindent
\emph{Addition} (after the Gille--Polo re-edition (2011), Exposé~I, 5.2.0 to
5.2.0.3; our formulation). The construction
\(\mathbf P\mapsto E(\mathbf P)=\underline{\operatorname{Hom}}(\mathbf G,\mathbf P)\)
is a functor from \((\mathbf O\text{-Mod.})\) to
\((\mathbf G\text{-}\mathbf O\text{-Mod.})\), and evaluation at the unit,
\(\varepsilon_{\mathbf P}(f)=f(1)\), is a morphism of \(\mathbf O\)-modules
\(E(\mathbf P)\to\mathbf P\). It makes \(E\) right adjoint to the forgetful
functor: for a \(\mathbf G\)-\(\mathbf O\)-module \(\mathbf M\),
\(\phi\mapsto\varepsilon_{\mathbf P}\circ\phi\) is a bijection from
\(\operatorname{Hom}_{\mathbf G\text{-}\mathbf O}(\mathbf M,E(\mathbf P))\)
onto \(\operatorname{Hom}_{\mathbf O}(\mathbf M,\mathbf P)\), whose inverse
sends \(u\) to the morphism \(m\mapsto(x\mapsto u(xm))\); indeed, if
\(\phi\) is equivariant, \(\phi(m)(x)=(x\phi(m))(1)=\phi(xm)(1)\). Hence
\(E\) carries injective \(\mathbf O\)-modules to injective
\(\mathbf G\)-\(\mathbf O\)-modules, the forgetful functor being exact. The
morphism corresponding to \(\operatorname{id}_{\mathbf F}\) is
\(\mu_{\mathbf F}\), and \(\varepsilon_{\mathbf F}\circ\mu_{\mathbf F}
=\operatorname{id}_{\mathbf F}\); so \(\mathbf F\) is a direct summand of
\(E(\mathbf F)\) as an \(\mathbf O\)-module.
\par\medskip

\medskip
\noindent\SGAInlineRefTarget{sga3:i:proposition:5-2-1}\textbf{Proposition 5.2.1.}
\textit{Suppose that \(\mathscr C\) is small, that finite products exist
in \(\mathscr C\), and that \(\mathbf G\) is representable.\footnote{Translator's note: the 1970 text has only ``Supposons
\(\mathbf G\) représentable.'' The Gille--Polo re-edition (2011),
Exposé~I, adds the other two hypotheses. Finite products are used in the
proof: they make \(\mathbf G^n\times\mathbf h_S\) representable, by
\(G^n\times S\), so that
\(\underline C^{\,n}(\mathbf G,\mathbf F)(S)=\mathbf F(G^n\times S)\) is
exact in \(\mathbf F\), exactness being tested at each object of
\(\mathscr C\); and \(\Gamma\) is then evaluation at the final object.
Smallness ensures that \((\mathbf G\text{-}\mathbf O\text{-Mod.})\) has
enough injective objects (3.2), so that derived functors are defined.
Without it, the proof shows that the \(\underline H^{\,n}(\mathbf G,-)\) and
\(H^n(\mathbf G,-)\) are the satellites of \(\underline H^{\,0}(\mathbf G,-)\)
and \(H^0(\mathbf G,-)\) in the sense of A.~Grothendieck, \textit{Sur quelques points d'algèbre homologique}, Tôhoku Math.\ J.\ \textbf{9} (1957), 119--221, 2.2.} Then the
functors \(\underline H^{\,n}(\mathbf G,-)\) (respectively,
\(H^n(\mathbf G,-)\)) are the derived functors of the left exact functor
\(\underline H^{\,0}(\mathbf G,-)\) (respectively,
\(H^0(\mathbf G,-)\)) on the category of
\(\mathbf G\)-\(\mathbf O\)-modules.}

By the well-known general results,\footnote{Translator's note: namely Grothendieck, \textit{loc.\ cit.}, 2.2.1 and 2.3
(references given in the Gille--Polo re-edition (2011), Exposé~I): an
effaceable cohomological functor is universal, and coincides with the
derived functors when there are enough injective objects.} it suffices to verify that the
\(\underline H^{\,n}(\mathbf G,-)\) (respectively, \(H^n(\mathbf G,-)\))
form a cohomological functor effaceable in dimensions \(>0\). Now
\(\underline C^{\,n}(\mathbf G,-)\), regarded as a functor on
\((\mathbf G\text{-}\mathbf O\text{-Mod.})\) with values in the category of
complexes of \((\mathbf O\text{-Mod.})\), is exact. This shows that the
\(\underline H^{\,n}(\mathbf G,-)\) indeed form a cohomological functor.
Since the functor \(\Gamma\) is exact, the same holds for the
\(H^n(\mathbf G,-)\). It will now suffice to prove the following.

\medskip
\noindent\SGAInlineRefTarget{sga3:i:lemma:5-2-2}\textbf{Lemma 5.2.2.}
\textit{For every \(\mathbf P\in\operatorname{Ob}(\mathbf O\text{-Mod.})\),
one has}
\[
  \underline H^{\,n}\bigl(
    \mathbf G,\underline{\operatorname{Hom}}(\mathbf G,\mathbf P)
  \bigr)=0,
  \qquad
  H^n\bigl(
    \mathbf G,\underline{\operatorname{Hom}}(\mathbf G,\mathbf P)
  \bigr)=0,
  \qquad n>0.
\]

It is enough to show that
\(\underline C^{\,*}
  (\mathbf G,\underline{\operatorname{Hom}}(\mathbf G,\mathbf P))\) and
\(C^*(\mathbf G,\underline{\operatorname{Hom}}(\mathbf G,\mathbf P))\)
are homotopically trivial in degrees \(>0\). It is enough, in fact, to do
this for the second complex, since the corresponding result for the first
follows by base change. Now there is the following homotopy operator:
\[
  (sf)(g_1,g_2,\ldots,g_n)(g)=f(g,g_1,\ldots,g_n)(e),
\]
where \(e\) denotes the unit element of \(\mathbf G(S)\) and \(f\) an
\((n+1)\)-cochain of \(\mathbf G(S)\) with values in
\(\underline{\operatorname{Hom}}(\mathbf G,\mathbf P)(S)\).\footnote{Translator's note: the 1970 text has \(f(g,g_0,g_1,\ldots,g_n)(e)\),
with one argument too many; the Gille--Polo re-edition (2011), Exposé~I,
writes the operator as above and checks that \(\partial s+s\partial\) is the
identity on cochains of degree \(n+1\geq1\). In our formulation: evaluate
\((\partial(sf)+s(\partial f))(g_1,\ldots,g_{n+1})\) at \(g\). The first term
of \(\partial(sf)\), \(g_1\cdot sf(g_2,\ldots,g_{n+1})\), gives
\(f(gg_1,g_2,\ldots,g_{n+1})(e)\), which cancels the term of \(s(\partial f)\)
in which \(g\) and \(g_1\) are multiplied; the other terms of
\(\partial(sf)\) cancel those of \(s(\partial f)\), whose signs are opposite
because \(g\) occupies the first place; what remains is the first term of
\(\partial f\), \((g\cdot f(g_1,\ldots,g_{n+1}))(e)=f(g_1,\ldots,g_{n+1})(g)\).}

\subsection*{\SGAInlineRefTarget{sga3:i:subsection:5-3}5.3. Cohomology of \(G\)-\(\mathscr O_S\)-modules}

Let \(S\) be a scheme, \(G\) an \(S\)-group affine over \(S\), and
\(\mathscr F\) a quasi-coherent \(G\)-\(\mathscr O_S\)-module. The cohomology groups of
\(G\) with values in \(\mathscr F\) are defined by
\[
  H^n(G,\mathscr F)
    =H^n\bigl(\mathbf h_G,\mathbf W(\mathscr F)\bigr)
\]
(for the notation, see \SGARefLink{sga3:i:subsection:4-6}{4.6}).

In view of \SGARefLink{sga3:i:proposition:4-6-4}{Proposition~4.6.4},
this cohomology is computed as follows: \(H^n(G,\mathscr F)\) is the
\(n\)-th cohomology group of the complex \(C^*(G,\mathscr F)\) defined
as follows:
\[
  C^n(G,\mathscr F)=\Gamma\bigl(S,\underline C^{\,n}(G,\mathscr F)\bigr),
  \qquad
  \underline C^{\,n}(G,\mathscr F)
    =\mathscr F\otimes
       \underbrace{\mathscr A(G)\otimes\cdots\otimes\mathscr A(G)}_
                   {n\text{ times}}.
\]
If \(f\) (respectively, \(a_i\)) is a section of \(\mathscr F\)
(respectively, of \(\mathscr A(G)\)) over an open subset of \(S\), then
\begin{align*}
  \partial(f\otimes a_1\otimes\cdots\otimes a_n)
    ={}&\mu_{\mathscr F}(f)\otimes a_1\otimes a_2\otimes\cdots\otimes a_n\\
     &+\sum_{i=1}^{n}(-1)^i
       f\otimes a_1\otimes\cdots\otimes\Delta a_i\otimes\cdots\otimes a_n\\
     &+(-1)^{n+1}
       f\otimes a_1\otimes a_2\otimes\cdots\otimes a_n\otimes1,
\end{align*}
where
\[
  \mu_{\mathscr F}:\mathscr F\longrightarrow
    \mathscr F\otimes\mathscr A(G),
  \qquad
  \Delta:\mathscr A(G)\longrightarrow\mathscr A(G)\otimes\mathscr A(G)
\]
describe the comodule structure of \(\mathscr F\). Notice in passing that the cohomology of
\(G\) with values in \(\mathscr F\) therefore depends only on the
comodule structure of \(\mathscr F\), and in particular only on the
\(S\)-monoid structure of \(G\).

In particular,
\[
  H^0(G,\mathscr F)=\Gamma(S,\mathscr F^G),
\]
where \(\mathscr F^G\), the \emph{sheaf of invariants of \(\mathscr F\)},
is the sheaf whose sections over an open subset \(U\) of \(S\) are those
sections of \(\mathscr F\) over \(U\) whose inverse image in every
\(S'\) over \(U\) is invariant under \(G(S')\).\footnote{Translator's note: the Gille--Polo re-edition (2011), Exposé~I, gives this
definition in 4.7, with an example showing that the sheaf
\(\mathscr F^G\) can be smaller than the functor of invariants: the natural
morphism \(\mathbf W(\mathscr F^G)\to\mathbf W(\mathscr F)^G\) is not an
isomorphism in general. In our formulation: let \(S=\operatorname{Spec}(\mathbb Z)\),
let \(G\) be the constant group \(\mathbb Z/2\mathbb Z\), and let its
generator act on \(\mathscr F=\mathcal O_S\) by \(-1\). A section \(x\) of
\(\mathcal O_S\) over an open set with \(-x=x\) is zero, the rings of
sections being torsion-free, so \(\mathscr F^G=0\); but for an
\(\mathbb F_2\)-algebra \(R\), every element of \(R\) is invariant, so
\(\mathbf W(\mathscr F)^G(\operatorname{Spec}R)=R\).}

\medskip
\noindent\SGAInlineRefTarget{sga3:i:theorem:5-3-1}\textbf{Theorem 5.3.1.}
\textit{Let \(S\) be an affine scheme and \(G\) an \(S\)-group that is
affine and flat over \(S\). The functors \(H^n(G,-)\) are the derived
functors of \(H^0(G,-)\) on the category of quasi-coherent
\(G\)-\(\mathscr O_S\)-modules.}

Since \(S\) is affine, and since \(G\) is affine and flat, one sees at once
that \(C^*(G,-)\) is indeed an exact functor on the preceding category.

It remains only to construct, for every \(\mathscr F\), a monomorphism of
\(\mathscr F\) into a \(G\)-\(\mathscr O_S\)-module \(E(\mathscr F)\) such
that \(H^n(G,E(\mathscr F))=0\) for \(n>0\). Now consider
\(E(\mathscr F)=\mathscr F\otimes\mathscr A(G)\), with the
\(G\)-\(\mathscr O_S\)-module structure defined by that of
\(\mathscr A(G)\), and the canonical homomorphism
\(\mu_{\mathscr F}:\mathscr F\to E(\mathscr F)\); axiom (CM~1) of
\SGARefLink{sga3:i:paragraph:4-7-2}{4.7.2} says exactly that \(\mu_{\mathscr F}\) is a morphism of
\(G\)-\(\mathscr O_S\)-modules. It now suffices to remark that the image of
\(\mu_{\mathscr F}\) under the functor \(\mathbf W\) is none other than the
effacement
\(\mu_{\mathbf W(\mathscr F)}:\mathbf W(\mathscr F)\to E(\mathbf W(\mathscr F))\)
defined in \SGARefLink{sga3:i:subsection:5-2}{5.2} (compare the corollary to
\SGARefLink{sga3:i:proposition:4-6-4}{Proposition~4.6.4}).\footnote{Translator's note: the 1970 text says ``le corollaire 1 à la
proposition 4.6.4''; Proposition 4.6.4 has a single corollary.}

\par\medskip\noindent
\emph{Addition} (after the Gille--Polo re-edition (2011), Exposé~I, 5.3.1.1
to 5.3.1.3; our formulation). We complete the proof. As \(G\) is flat over
\(S\), the quasi-coherent \(G\)-\(\mathscr O_S\)-modules form an abelian
category, exactness being tested on the underlying \(\mathscr O_S\)-modules
(Addition after 4.7.2); the functors
\(\mathscr F\mapsto\mathscr F\otimes\mathscr A(G)^{\otimes n}\) are exact,
and so is \(\Gamma(S,-)\) on quasi-coherent modules, \(S\) being affine; this
gives the exactness of \(C^*(G,-)\). For a quasi-coherent
\(\mathscr O_S\)-module \(\mathscr P\), let \(\operatorname{Ind}(\mathscr P)\)
be \(\mathscr P\otimes\mathscr A(G)\) with the coaction
\(\operatorname{Id}\otimes\Delta\). For every comodule \(\mathscr M\),
\(\phi\mapsto(\operatorname{Id}\otimes\varepsilon)\circ\phi\) is a bijection
from comodule morphisms \(\mathscr M\to\operatorname{Ind}(\mathscr P)\) onto
\(\mathscr O_S\)-linear maps \(\mathscr M\to\mathscr P\), with inverse
\(u\mapsto(u\otimes\operatorname{Id})\circ\mu_{\mathscr M}\), as (CM~1),
(CM~2) and the counit axiom of \(\mathscr A(G)\) show. So
\(\operatorname{Ind}\) is right adjoint to the exact forgetful functor and
carries injective objects to injective objects. The module
\(E(\mathscr F)\) above is \(\operatorname{Ind}(\mathscr F)\); by the
corollary to 4.6.4,
\(\mathbf W(\operatorname{Ind}(\mathscr P))\simeq
\underline{\operatorname{Hom}}(\mathbf h_G,\mathbf W(\mathscr P))\), the
coaction \(\operatorname{Id}\otimes\Delta\) corresponding to the action
\((gf)(x)=f(xg)\) of 5.2, so \(H^n(G,\operatorname{Ind}(\mathscr P))=0\) for
\(n>0\) by 5.2.2. Finally, \(S\) being affine, every quasi-coherent
\(\mathscr F\) embeds in an injective object \(\mathscr I\) of the
quasi-coherent modules; then
\(\mathscr F\to\operatorname{Ind}(\mathscr F)\to\operatorname{Ind}(\mathscr I)\)
is a monomorphism of comodules (the second map because \(\mathscr A(G)\) is
flat) into an injective object. Hence the category has enough injective
objects, and the \(H^n(G,-)\) are the derived functors of \(H^0(G,-)\)
(A.~Grothendieck, Tôhoku Math.\ J.\ \textbf{9} (1957), 2.2.1 and 2.3).
Compare J.~C.~Jantzen,
\textit{Representations of algebraic groups}, 2nd ed., American
Mathematical Society, 2003, Part~I, 3.3--3.4, 3.9, 4.2 and 4.14--4.16,
where group schemes are affine (\textit{loc.\ cit.}, 2.1).
\par\medskip

Note now that axiom (CM~2) shows that, as an \(\mathscr O_S\)-module,
\(\mathscr F\) is a direct summand of \(E(\mathscr F)\). This implies:

\medskip
\noindent\SGAInlineRefTarget{sga3:i:proposition:5-3-2}\textbf{Proposition 5.3.2.}
\textit{Let \(S\) be an affine scheme and \(G\) an \(S\)-group that is
affine and flat. Suppose that every exact sequence
\[
  0\longrightarrow\mathscr F_1
   \longrightarrow\mathscr F_2
   \longrightarrow\mathscr F_3
   \longrightarrow0
\]
of quasi-coherent \(G\)-\(\mathscr O_S\)-modules which splits as a sequence
of \(\mathscr O_S\)-modules also splits as a sequence of
\(G\)-\(\mathscr O_S\)-modules. The functors \(H^n(G,-)\), \(n>0\), vanish
(or, equivalently, the functor \(H^0(G,-)\) is exact).}

Indeed, by hypothesis, the sequence
\[
  0\longrightarrow\mathscr F
   \longrightarrow E(\mathscr F)
   \longrightarrow E(\mathscr F)/\mathscr F
   \longrightarrow0
\]
is split. Thus \(\mathscr F\) is a direct summand of \(E(\mathscr F)\),
while the cohomology of the latter vanishes.

Combining this with \SGARefLink{sga3:i:proposition:4-7-4}{Proposition~4.7.4} immediately gives:

\medskip
\noindent\SGAInlineRefTarget{sga3:i:theorem:5-3-3}\textbf{Theorem 5.3.3.}
\textit{Let \(S\) be an affine scheme and \(G\) a diagonalizable
\(S\)-group. For every quasi-coherent \(G\)-\(\mathscr O_S\)-module
\(\mathscr F\), one has}
\[
  H^n(G,\mathscr F)=0\qquad(n>0).
\]

\medskip
\noindent\textbf{Remark.}
\SGARefLink{sga3:i:proposition:5-3-2}{Proposition~5.3.2} remains valid when \(G\) is not necessarily flat over
\(S\); the proof then uses relative cohomology.

\par\bigskip\noindent\SGAInlineRefTarget{sga3:i:section:6}%
\emph{Addition} (after the Gille--Polo re-edition (2011), Exposé~I, \S6,
``Objets et modules \(G\)-équivariants''; our formulation). The re-edition
ends this Exposé with a section on equivariant objects and modules, used in
Exposés~II, III and VI\textsubscript{B}. We give its content in our own
terms; the numbers in brackets are those of the re-edition.

\par\smallskip\noindent\textit{(a) Equivariant objects} [6.1, 6.2].
Let \(\mathbf G\) be a \(\widehat{\mathscr C}\)-group acting on an object
\(X\) of \(\widehat{\mathscr C}\) by \(\lambda:\mathbf G\times X\to X\), and
let \(\pi:M\to X\) be a morphism of \(\widehat{\mathscr C}\). A
\emph{\(\mathbf G\)-equivariant structure} on \(M\) over \(X\) is an action
\(\Lambda:\mathbf G\times M\to M\) with
\(\pi\circ\Lambda=\lambda\circ(\operatorname{id}\times\pi)\). For
\(U\in\operatorname{Ob}(\mathscr C)\) and \(x\in X(U)\), let
\(M_x=\mathbf h_U\times_{x,X}M\) be the part of \(M\) above \(x\); for
\(g\in\mathbf G(U)\), \(\Lambda\) restricts to morphisms
\(\Lambda_x(g):M_x\to M_{gx}\) over \(U\), compatible with base change, with
\(\Lambda_x(1)=\operatorname{id}\) and
\(\Lambda_{hx}(g)\circ\Lambda_x(h)=\Lambda_x(gh)\); so each \(\Lambda_x(g)\)
is an isomorphism, with inverse \(\Lambda_{gx}(g^{-1})\), and such data
determine \(\Lambda\). If \(\mathbf O\) is a \(\widehat{\mathscr C}\)-ring and
\(M\) an \(\mathbf O_X\)-module (Addition in 4.3), \(M\) is a
\emph{\(\mathbf G\)-equivariant \(\mathbf O_X\)-module} if moreover the
\(\Lambda_x(g)\) are \(\mathbf O\)-linear. In one morphism: the
\(\Lambda_x(g)\) amount to a morphism of
\(\mathbf O_{\mathbf G\times X}\)-modules
\(\theta:(\mathbf G\times X)\times_{\mathrm{pr}_X}M\to(\mathbf G\times X)\times_\lambda M\),
\((g,x,m)\mapsto(g,x,gm)\), and \(\Lambda\) is an action if and only if
\(\theta\) is the identity above the unit section and satisfies the cocycle
condition
\((\pi_{\mathbf G}\times\operatorname{id}_X)^*\theta
=(\operatorname{id}_{\mathbf G}\times\lambda)^*\theta\circ\mathrm{pr}_{23}^*\theta\)
over \(\mathbf G\times\mathbf G\times X\), both sides going from the
pull-back of \(M\) by \((g,h,x)\mapsto x\) to its pull-back by
\((g,h,x)\mapsto ghx\) (at a point, \(g(hm)=(gh)m\)). If \(\theta\) is an
isomorphism, the unit condition follows from the cocycle condition; for a
monoid \(\mathbf G\) it must be imposed.

\par\smallskip\noindent\textit{(b) Equivariant morphisms, stabilizers} [6.3].
If \(\mathbf G\) also acts on \(Y\), a morphism \(f:Y\to X\) is equivariant
if \(f\circ\lambda_Y=\lambda_X\circ(\operatorname{id}\times f)\); then
\(f^*M=Y\times_{f,X}M\), with \(g(y,m)=(gy,gm)\), is a
\(\mathbf G\)-equivariant \(\mathbf O_Y\)-module [6.3.1]. A morphism of
\(\mathbf O_X\)-modules \(u:M\to N\) is equivariant if
\(u\circ\Lambda_M=\Lambda_N\circ(\operatorname{id}\times u)\). For an
arbitrary \(f:Y\to X\), \(\mathbf G\) acts on
\(\underline{\operatorname{Hom}}(Y,X)\) by
\(g\cdot\phi=\lambda_X(g)\circ\phi\circ\lambda_Y(g)^{-1}\); the stabilizer
\(\underline{\operatorname{Stab}}_{\mathbf G}(f)\) of the section \(f\)
(Addition in 2.3) consists of the \(g\) with
\(\lambda_X(g)\circ f_T=f_T\circ\lambda_Y(g)\), and \(f\) is equivariant for
it [6.3.2].

\par\smallskip\noindent\textit{(c) Sections} [6.4].
Let \(\mathscr C\) have a final object \(S_0\), and let \(M\) be a
\(\mathbf G\)-equivariant \(\mathbf O_X\)-module. The functor of sections
\(\prod_{X/S_0}M\) (Exposé~II, \SGARefLink{sga3:ii:definition:1-1}{1.1}) assigns to \(T\) the
\(\mathbf O(X\times T)\)-module \(\operatorname{Hom}_X(X\times T,M)\). For
\(g\in\mathbf G(T)\) and a section \(s\), put \((gs)(x,t)=g\,s(g^{-1}x,t)\)
at the points \((x,t)\) of \(X\times T\) (with \(g\) evaluated at \(t\)).
This is again a section, \(\mathbf G(T)\) acts in this way, functorially in
\(T\), and \(g(\alpha s)=(g\alpha)(gs)\) for
\(\alpha\in\mathbf O(X\times T)\), where \((g\alpha)(x,t)=\alpha(g^{-1}x,t)\)
[6.4.1]. The functions coming from \(\mathbf O(T)\) are fixed; so
\(\prod_{X/S_0}M\) is a \(\mathbf G\)-\(\mathbf O\)-module. For any
\(\phi:Y\to X\), \(\phi^*M\) is equivariant for
\(\underline{\operatorname{Stab}}_{\mathbf G}(\phi)\), by (b), whence a
\(\underline{\operatorname{Stab}}_{\mathbf G}(\phi)\)-\(\mathbf O\)-module
\(\prod_{Y/S_0}\phi^*M\) [6.4.2].

\par\smallskip\noindent\textit{(d) Equivariant sheaves} [6.5].
Let \(G\) be an \(S\)-group scheme acting on an \(S\)-scheme \(X\) by
\(\lambda\), and \(\mathscr F\) an \(\mathcal O_X\)-module, not necessarily
quasi-coherent. One calls \(\mathscr F\) \emph{\(G\)-equivariant} if
\(\mathbf W(\mathscr F)\), the \(\mathbf O_X\)-module
\((x:U\to X)\mapsto\Gamma(U,x^*\mathscr F)\), is \(G\)-equivariant [6.5.1].
As \(\mathbf W\) is fully faithful (4.6.2 and the Note in 4.6.1), this
amounts to isomorphisms \(x^*\mathscr F\simeq(gx)^*\mathscr F\) for the
points \((g,x)\) of \(G\times_SX\), compatible with base change and
satisfying the cocycle relation; by taking the universal point, to an
isomorphism of \(\mathcal O_{G\times_SX}\)-modules
\(\theta:\mathrm{pr}_X^*\mathscr F\xrightarrow{\sim}\lambda^*\mathscr F\)
satisfying the cocycle condition of (a) on \(G\times_SG\times_SX\). The
equivariant morphisms \(u:\mathscr E\to\mathscr F\) are those with
\(\theta_{\mathscr F}\circ\mathrm{pr}_X^*(u)=\lambda^*(u)\circ\theta_{\mathscr E}\),
and pull-backs by equivariant morphisms are equivariant [6.5.2]. If
\(\mathscr F\) is quasi-coherent, \(\theta^{-1}\) induces, through the
contravariant functor \(\mathbb V\) (4.6.3), an isomorphism
\((G\times_SX)\times_{\mathrm{pr}_X}\mathbb V(\mathscr F)\to(G\times_SX)\times_\lambda\mathbb V(\mathscr F)\),
whose composite with the second projection is an action of \(G\) on
\(\mathbb V(\mathscr F)\) above that on \(X\), linear on the fibres;
equivalently, \(\mathbf V(\mathscr F)\) is a \(G\)-equivariant
\(\mathbf O_X\)-module. For invertible \(\mathscr F\) this is a
\(G\)-linearization in the sense of D.~Mumford, \textit{Geometric invariant
theory}, Springer, 1965, Chapter~1, \S3, Definition~1.6, where the
isomorphism is written in the other direction [6.5.3].
\SGAInlineRefTarget{sga3:i:remark:6-5-4}When \(X=S\) with the trivial action,
a \(G\)-equivariant \(\mathcal O_S\)-module is a
\(G\)-\(\mathcal O_S\)-module (4.7.1): if \(f:G\to S\) is the structural
morphism, \(\theta\in\operatorname{Aut}(f^*\mathscr F)
=\underline{\operatorname{Aut}}_{\mathbf O_S}(\mathbf W(\mathscr F))(G)\) is
the image of the universal point under \(\rho\), and corresponds by
adjunction to \(\mathscr F\to f_*f^*\mathscr F\), which for \(G\) affine is
the comodule map of 4.7.2 [6.5.4, used in Exposé~VI\textsubscript{B},
11.10~bis].

\par\smallskip\noindent\textit{(e) Direct images} [6.6].
Let \(G\) be flat over \(S\) and act on \(X\) and \(Y\), let \(p:X\to Y\) be
equivariant, quasi-compact and quasi-separated, and let \(\mathscr F\) be a
quasi-coherent \(G\)-equivariant \(\mathcal O_X\)-module. Then
\(p_*\mathscr F\) is quasi-coherent (EGA~I, 9.2.1) and \(G\)-equivariant
[6.6.1]. Indeed, put \(q=\operatorname{id}_G\times p\); the squares formed by
\(q\), \(p\) and the pairs \((\mathrm{pr}_X,\mathrm{pr}_Y)\),
\((\lambda_X,\lambda_Y)\) are cartesian (for the second,
\((g,x)\mapsto(g,p(x),gx)\) identifies \(G\times_SX\) with
\((G\times_SY)\times_{\lambda_Y,Y,p}X\)), and \(\mathrm{pr}_Y\) and
\(\lambda_Y\) are flat, \(\lambda_Y\) being \(\mathrm{pr}_Y\) composed with
the automorphism \((g,y)\mapsto(g,gy)\). By flat base change (EGA~III,
1.4.15, and EGA~IV\textsubscript{1}, 1.7.21),
\(q_*\mathrm{pr}_X^*\mathscr F=\mathrm{pr}_Y^*p_*\mathscr F\) and
\(q_*\lambda_X^*\mathscr F=\lambda_Y^*p_*\mathscr F\), so \(q_*\theta\) is an
isomorphism \(\mathrm{pr}_Y^*p_*\mathscr F\to\lambda_Y^*p_*\mathscr F\), and
its cocycle condition follows in the same way over \(G\times_SG\times_SY\).
For \(Y=S\) with the trivial action, \(p_*\mathscr F\) is therefore a
quasi-coherent \(G\)-\(\mathcal O_S\)-module, an \(\mathscr A(G)\)-comodule
if \(G\) is affine (4.7.2; see also Exposé~VI\textsubscript{B}, \SGARefLink{concl:sga3-VIB-11.6.1-ii}{11.6.1~(ii)},
for \(G\) quasi-compact and quasi-separated and \(p_*\mathscr F\) flat), and
the canonical morphism from \(\mathbf W(p_*\mathscr F)\) to the
\(G\)-\(\mathbf O_S\)-module \(\prod_{X/S}\mathbf W(\mathscr F)\) of (c), given
at \(f:T\to S\) by
\(\Gamma(T,f^*p_*\mathscr F)\to\Gamma(X_T,f_X^*\mathscr F)\), is compatible
with the actions of \(G\), and bijective for \(T\) flat over \(S\) [6.6.2;
compare Mumford, \textit{loc.\ cit.}, p.~32].

\par\smallskip\noindent\textit{(f) Stabilizers} [6.7].
Let \(G\) act on \(X\) and \(Y\), let \(\phi:Y\to X\) be an \(S\)-morphism,
not necessarily equivariant, and let \(\mathscr F\) be a quasi-coherent
\(G\)-equivariant \(\mathcal O_X\)-module. Then
\(\underline{\operatorname{Stab}}_G(\phi)=\prod_{P/G}P'\), where
\(P=G\times_SY\), \(P'\) is the inverse image of the diagonal of \(X\) under
\((g,y)\mapsto(\phi(y),g\phi(g^{-1}y))\), and the Weil restriction is along
the projection \(P\to G\). So, if \(X\) is separated over \(S\) and \(Y\) is
essentially free over \(S\), the stabilizer is a closed subgroup scheme
\(H\) of \(G\) (Exposé~VI\textsubscript{B}, \SGARefLink{def:sga3-VIB-6.2.1}{6.2.1} and \SGARefLink{ex:sga3-VIB-6.2.4-a}{6.2.4~(a)}); then
\(\phi^*\mathscr F\) is a quasi-coherent \(H\)-equivariant
\(\mathcal O_Y\)-module, by (b), and if \(H\) is flat over \(S\) and
\(Y\to S\) is quasi-compact and quasi-separated, its direct image on \(S\)
is an \(H\)-\(\mathcal O_S\)-module, by (e). (The re-edition states this last
point with \(G\), instead of \(H\), flat over \(S\); (e) needs the flatness
of the group that acts.) For a section \(\tau:S\to X\) of a separated
\(X\), the stabilizer is the closed subgroup scheme \(H=G\times_XS\), fibre
product of \(g\mapsto g\tau\) and \(\tau\), and \(\tau^*\mathscr F\) is a
quasi-coherent \(H\)-\(\mathcal O_S\)-module, by (b) and (d); it is an
\(\mathscr A(H)\)-comodule if \(H\) is affine over \(S\), for instance if
\(G\) is [6.7.1].

\par\smallskip\noindent\SGAInlineRefTarget{sga3:i:subsection:6-8}\textit{(g) Equivariant sheaves on a group} [6.8].
Let \(\pi:G\to S\) be an \(S\)-group scheme with unit section
\(\varepsilon\), acting on itself by left translations.
\SGAInlineRefTarget{sga3:i:proposition:6-8-1}The functors
\(\mathscr E\mapsto\pi^*\mathscr E\) and
\(\mathscr F\mapsto\varepsilon^*\mathscr F\) are quasi-inverse equivalences
between the quasi-coherent \(\mathcal O_S\)-modules and the quasi-coherent
\(G\)-equivariant \(\mathcal O_G\)-modules [6.8.1]. Indeed, as
\(\pi\circ\pi_G=\pi\circ\mathrm{pr}_2\) on \(G\times_SG\), the identity
\(\mathrm{pr}_2^*\pi^*\mathscr E=\pi_G^*\pi^*\mathscr E\) is a \(\theta\)
satisfying the cocycle condition. Let
\(\sigma=(\operatorname{id}_G,\varepsilon\pi):G\to G\times_SG\), so that
\(\mathrm{pr}_2\circ\sigma=\varepsilon\pi\) and
\(\pi_G\circ\sigma=\operatorname{id}_G\). For an equivariant \(\mathscr F\),
\(\sigma^*\theta\) is an isomorphism
\(\pi^*\varepsilon^*\mathscr F\to\mathscr F\), equivariant by the cocycle
condition at the points \((g,h,e)\). For an equivariant morphism
\(u:\pi^*\mathscr E\to\pi^*\mathscr E'\), applying \(\sigma^*\) to the
condition of (d) gives \(u=\pi^*\varepsilon^*u\); with
\(\varepsilon^*\pi^*=\operatorname{id}\), this proves the assertion.

\SGAInlineRefTarget{sga3:i:remark:6-8-2}Let \(G\times_SG\) act on \(G\) by
\((g_1,g_2)\cdot g=g_1gg_2^{-1}\). The stabilizer of \(\varepsilon\) is the
diagonal \(H\simeq G\); so, for a quasi-coherent
\((G\times_SG)\)-equivariant \(\mathcal O_G\)-module \(\mathscr F\),
\(\varepsilon^*\mathscr F\) is an \(H\)-\(\mathcal O_S\)-module, by (f), and
\(\mathscr F\mapsto\varepsilon^*\mathscr F\) is an equivalence onto the
quasi-coherent \(G\)-\(\mathcal O_S\)-modules [6.8.2]. The re-edition
derives this from descent (Exposé~IV, \S2; SGA~1, Exposé~VIII;
R.~W.~Thomason, \textit{Equivariant resolution, linearization, and
Hilbert's fourteenth problem over arbitrary base schemes}, Adv.\ Math.\
\textbf{65} (1987), 16--34, 1.2--1.3). It also follows from the preceding
equivalence: restricted to \(G\times1\), \(\mathscr F\) becomes
\(\pi^*\mathscr E\), \(\mathscr E=\varepsilon^*\mathscr F\), with
\(G\times1\) acting through identity maps of \(\mathscr E\); since
\((1,b)=(gb^{-1}g^{-1},1)\,(gbg^{-1},b)\) and \((gbg^{-1},b)\) is conjugate
by \((g,1)\) to the diagonal element \((b,b)\), the cocycle condition shows
that \((g_1,g_2)\) acts at every point through the action of
\(g_2\in H\) on \(\mathscr E\); conversely, this rule defines a
\((G\times_SG)\)-equivariant structure on \(\pi^*\mathscr E\) for every
\(G\)-\(\mathcal O_S\)-module \(\mathscr E\).

\SGAInlineRefTarget{sga3:i:remark:6-8-3}For a quasi-coherent
\(\mathcal O_S\)-module \(\mathscr E\), the sections of
\(\pi_*\pi^*\mathscr E\) over an open \(V\) of \(S\) that are invariant
under \(G\), those \(\gamma\) with \(g\gamma_{S'}=\gamma_{S'}\) for all
\(S'\to V\) and \(g\in G(S')\), are exactly the \(\pi^*e\),
\(e\in\Gamma(V,\mathscr E)\): invariance gives
\(a^*\gamma=\mathrm{pr}_2^*\gamma\) over \(G_V\times_VG_V\), where
\(a(g,h)=g^{-1}h\), and pulling back by the diagonal gives
\(\gamma=\pi^*\varepsilon^*\gamma\). So
\(\mathscr E\to\pi_*^G\pi^*\mathscr E\) is an isomorphism, without
further hypothesis [6.8.3].

If \(H\) is a flat \(S\)-group scheme acting on an \(S\)-scheme \(X\) flat
over \(S\), then \(H\times_SX\) and \(H\times_SH\times_SX\) are flat over
\(S\), so an \(H\)-equivariant structure on an \(\mathcal O_X\)-module
\(\mathscr F\) (\(\theta\) and its cocycle condition in (d)) may be given,
and checked, on
the restriction of \(\mathbf W(\mathscr F)\) to points with values in
\(S\)-schemes flat over \(S\) [6.8.4; the re-edition does not require \(X\)
to be flat, but the universal points must be among these test schemes].
Let \(X\) be an \(S\)-group scheme and \(Y\) a subgroup scheme of \(X\),
flat over \(S\), with conormal sheaf \(\mathscr N_{Y/X}\) (EGA~IV\textsubscript{4},
16.1.2). Then \(\mathscr N_{Y/X}\) is a \((Y\times_SY)\)-equivariant
\(\mathcal O_Y\)-module, \(Y\times_SY\) acting on \(Y\) by
\((y_1,y_2)\cdot y=y_1yy_2^{-1}\) [6.8.6]. Indeed, the conormal sheaf
commutes with flat base change \(S'\to S\) (\textit{loc.\ cit.},
16.2.2~(iii)) [6.8.5]; for \(S'\) flat over \(S\), a point \((y_1,y_2)\) of
\(Y\times_SY\) with values in \(S'\) acts on \(X_{S'}\) by the automorphism
\(x\mapsto y_1xy_2^{-1}\), which preserves \(Y_{S'}\) and so induces
isomorphisms between the pull-backs of
\(\mathscr N_{Y_{S'}/X_{S'}}=\mathscr N_{Y/X}\otimes\mathcal O_{Y_{S'}}\) at a
point \(y\) and at its transform; these satisfy the cocycle relation, and
the preceding remark applies, \(Y\) and \(Y\times_SY\) being flat over
\(S\).
\par\medskip

